\section{Grammaire}

    \subsection{Les particules}

        \subsubsection{La particule の}

    \begin{itemize}
        \item \textbf{La possession :} La particule の indique la possession entre deux éléments, par la formule 
        \begin{center}
            [possesseur の possédé]
        \end{center}
        \begin{tcolorbox}
            -- 山本さんのお母さんは高校の先生です。

            -- やまもとさん　の　おかあさん　は　こうこう　の　せんせい　です。

            -- La mère de M. Yamamoto est professeure au lycée.
        \end{tcolorbox}
    \end{itemize}

        \subsubsection{La particule と}

    \begin{itemize}
        \item \textbf{L'énumération exhaustive :} Cette particule se place entre deux \textit{noms}, pour les énumérer, par la formule 
        \begin{center}
            [nom1 と nom2]
        \end{center}
        \item \textbf{« Avec » :} En la plaçant après le nom d'un \textit{être vivant}, 
        \begin{center}
            [nom と],
        \end{center}
         signifie « avec \textit{nom} ».
        \begin{tcolorbox}
            -- 土曜日に、友達といっしょに図書館で日本語の本を読みました。

            -- どようび　に、ともだち　と　いっしょに　としょかん　で　にほんご　の　ほん　を　よみますした。

            -- Samedi, j'ai lu un livre en japonais à la bibliothèque avec un ami.
        \end{tcolorbox}

        \item \textbf{Citation et pensée :} On peut utiliser と pour encadrer une expression ou phrase que l'on veut citer. Plus généralement, on peut la placer juste avant des verbes comme 言う et 思う.
        
        \begin{tcolorbox}
            -- 明日は雨が降ると思います。

            -- あしたは　あめが　ふる　と　おもいます。

            -- Je pense qu'il pleuvra demain.
        \end{tcolorbox}
    \end{itemize}

        \subsubsection{La particule や}

    \begin{itemize}
        \item \textbf{L'énumération non exhaustive :} Cette particule se place entre deux \textit{noms}, pour énumérer une liste non exhaustive, par la formule \begin{center}
            [nom1 や nom2]
        \end{center}
    \end{itemize}

        \subsubsection{La particule は}

    Cette particule, notée は, est prononcée « wa ».

    \begin{itemize}
        \item \textbf{Le thème :} Cette particule indique le thème de la phrase, a priori différent du sujet du verbe. On peut traduire l'écriture 
        \begin{center}
            [thème は]
        \end{center}
        par « en ce qui concerne \textit{thème} ».
    \end{itemize}

        \subsubsection{La particule に}

    \begin{itemize}
        \item \textbf{Le lieu (passif) :} Accompagné d'un verbe de présence, un lieu suivi de に 
        \begin{center}
            [lieu に \dots verbe-passif]
        \end{center}
        indique le lieu où il n'y a pas d'action.

        \begin{tcolorbox}
            -- 新聞は机にある。

            -- しんぶん　は　つくえ　に　ある。

            -- Le journal est sur le bureau.
        \end{tcolorbox}

        \item \textbf{Le point d'arrivée d'un déplacement :} S'il y a un déplacement, on peut indiquer le point d'arrivée avec la formule 
        \begin{center}
            [point d'arrivée に \dots verbe-déplacement]
        \end{center}
        \begin{tcolorbox}
            -- 今日学校に行きません。

            -- きょう　がっこうに　いきません。

            -- Je ne vais pas à l'école aujourd'hui.
        \end{tcolorbox}

        \item \textbf{Le COI :} La particule に indique le COI associé au verbe par la forme 
        \begin{center}
            [\dots　COI　に　verbe]
        \end{center}

        \item \textbf{Indiquer un temps absolu :} La particule に indique la temporalité absolue (telle heure, tel jour, tel mois) à laquelle une action va être réalisée par la forme 
        \begin{center}
            [date / heure　に　verbe-action]
        \end{center}
        On ne peut pas utiliser に accompagné des notions relatives de temps, comme aujourd'hui ou tous les jours.
        \begin{tcolorbox}
            -- 日曜日に京都に行く。
            
            -- にちよう　に　きょとう　に　いく。

            -- J'irai à Kyoto dimanche.
        \end{tcolorbox}
    \end{itemize}

        \subsubsection{La particule で}

    \begin{itemize}
        \item \textbf{Le lieu d'une action :} Accompagné d'un verbe d'action, un lieu suivi de で
        \begin{center}
            [lieu で \dots verbe-action]
        \end{center}
        indique le lieu dans lequel se déroule une action.
        \begin{tcolorbox}
            -- 図書館で本を読みます。

            -- としょかん　で　ほん　を　よみます。

            -- Je lis un livre dans une bibliothèque.
        \end{tcolorbox}

        \item \textbf{Le moyen, l'instrument et le transport :} La particule で permet de préciser le moyen par lequel est réalisé le verbe :
        \begin{center}
            [moyen で \dots verbe-action]
        \end{center}
        Il peut s'agir notamment d'un moyen de transport.

        \item \textbf{Coordonner un な-adjectifs avec un adjectif :} La particule で permet de coordonner un な-adjectif avec un autre adjectif :
        
        \begin{center}
            「な-adjectif　で　な-adjectif　な　nom」

            「な-adjectif　で　い-adjectif　nom」
        \end{center}

        \begin{tcolorbox}
            -- 旅行した町は静かできれいなところでしたが、少し不便でした。

            -- りょこうした　まち　は　しずか　で　きれい　な　ところ　でした　が、すこし　ふべん　でした。

            -- La ville où j'ai voyagé était calme et jolie, mais assez peu pratique.
        \end{tcolorbox}

        \item \textbf{La cause ou la raison simple :} Avec des noms représentant des événements ou des états, で permet d'en exprimer la cause.
        
        \begin{tcolorbox}
            -- 風邪で休みました。

            -- ふうじゃで　やすみました。

            -- J'ai été absent à cause d'un rhume.
        \end{tcolorbox}
    \end{itemize}

        \subsubsection{La particule へ}

    Cette particule, bien que notée へ, se prononce « e ».

    \begin{itemize}
        \item \textbf{La direction :} Accompagné d'un verbe de déplacement directionnel (avancer, aller, se diriger vers...), un lieu suivi de へ
        \begin{center}
            [lieu へ \dots verbe-directionnel]
        \end{center}
        indique la direction vers laquelle on se déplace. Elle s'utilise surtout avec les verbes 行く, 来る, et 帰る.
    \end{itemize}

        \subsubsection{La particule を}

    Cette particule, notée を, se prononce « o ».

    \begin{itemize}
        \item \textbf{Le COD :} La particule を permet d'indiquer le COD du verbe : 
        \begin{center}
            [\dots COD を verbe]
        \end{center}
    \end{itemize}

        \subsubsection{か}

    \begin{itemize}
        \item \textbf{La question :} La particule か en fin de phrase indique sa forme interrogative. 
        
        \begin{tcolorbox}
            -- 留学生ですか。田中さんは何歳ですか。

            -- りゅうがくせい　ですか。たなかさんは　なんさい　ですか。

            -- Es-tu un étudiant international ? Quel âge avez-vous M. Tanaka ?
        \end{tcolorbox}
       
        \item \textbf{L'invitation :} La particule か en fin de phrase, précédée d'un verbe à la forme négative permet de formuler une invitation.
        
        \begin{center}
            [\dots　verbe-négatif　か]
        \end{center}

        \begin{tcolorbox}
            -- 昼ご飯を食べませんか。

            -- ひるごはん　を　たべません　か。

            -- Voulez-vous manger avec moi ce midi ?
        \end{tcolorbox}

        \item \textbf{Le choix entre deux alternatives :} Répétée deux fois, la particule か permet de présenter un choix entre deux alternatives.
        
        \begin{tcolorbox}
            -- 今日か明日か分かりません。学生ですか、先生ですか

            -- きょうか　あしたか　わかりません。がくせい　ですか、せんせい　ですか。

            -- Je ne sais pas si c'est aujourd'hui ou demain.　Êtes-vous étudiant ou professeur ?
        \end{tcolorbox}
    \end{itemize}

        \subsubsection{も}

    \begin{itemize}
        \item \textbf{Aussi :} La particule も située après un groupe de mot permet d'exprimer l'idée de similarité avec la phrase évoquée avant.
        \begin{center}
            [groupe も \ldots]
        \end{center}
        \begin{tcolorbox}
            -- たけしさんは日本人です。陽子さんも日本人です。

            -- たけしさんはにほんじんです。ようこさんも日本人です。

            -- Takeshi est japonais. Yôko aussi.
        \end{tcolorbox}
        La particule も est très proche de la particule は, et s'utilise souvent aux mêmes endroits de la phrase. Elle indique aussi le thème, mais là où は introduit une notion d'exclusion (nouveau thème, différent), la particule も introduit une notion d'inclusion. Rarement avec des phrases positives, les locutions 何も (tout), 誰も (tout le monde)\dots, associés à de la négation, permettent d'exprimer l'idée de aucun, personne\dots

        \begin{tcolorbox}
            -- 誰もいない。

            -- だれも いない。

            -- Il n'y a personne.
        \end{tcolorbox}
    \end{itemize}

        \subsubsection{ね}

    \begin{itemize}
        \item \textbf{« N'est-ce pas ? » :} La particule ね, placée en fin de phrase, indique la recherche de confirmation de la part de l'autre.
        \begin{center}
            [\ldots ね]
        \end{center}
        \begin{tcolorbox}
            -- これは肉じゃないですね。

            -- これ　は　にく　じゃないです　ね。

            -- Ce n'est pas de la viande, si ?
        \end{tcolorbox}
    \end{itemize}

        \subsubsection{よ}

    \begin{itemize}
        \item \textbf{Accentuation} La particule よ, placée en fin de phrase, indique une forte croyance en ce que l'on dit (« Crois-moi »). Cette particule peut remplacer un point d'exclamation, et met de façon générale une emphase sur notre propos.
        \begin{center}
            [\ldots よ]
        \end{center}
    \end{itemize}

        \subsubsection{て}

    \begin{itemize}
        \item \textbf{Coordonner un い-adjectif avec un adjectif :} La particule て permet de coordonner un い-adjectif de radical <adj> avec un autre adjectif, par la forme
    
        \begin{center}
            「<adj> + く　て　な-adjectif　な　nom」

            「<adj> + く　て　い-adjectif　nom」
        \end{center}

        \begin{tcolorbox}
            -- 昨日見た映画は長くて面白かったです。

            -- きのう　みた　えいが　は　ながく　て　おもしろい　かったです。

            -- Le film que j'ai vu hier était long et intéressant.
        \end{tcolorbox}
    \end{itemize}

        \subsubsection{が}

    \begin{itemize}
        \item \textbf{Le sujet :} La particule が permet, dans le cas où le thème est différent du sujet du verbe, de l'indiquer.
        \begin{center}
            「\dots　sujet　+　が　\dots　verbe」
        \end{center}
        Voici différents contextes d'utilisation de cette particule sujet :
        \begin{itemize}
            \item \textit{Aimer / Détester :} Pour exprimer la notion de goût, on utilise la particule が associée aux な-adjectifs 好き (すき, aimer) ou 嫌い (きらい, détester). 
            \begin{tcolorbox}
                -- 私は日本の音楽が好きですが、弟はそれが嫌いです。

                -- わたし　は　にほん　の　おんがく　が　すき　です　が、おとうと　は　それ　が　きらい　です。

                -- J'aime la musique japonaise, mais mon petit frère déteste ça.
            \end{tcolorbox}
            \item \textit{Insister sur le sujet ou donner une nouvelle information :} On peut utiliser が pour insister sur le sujet du verbe, notamment lorsque l'on répond à une question, pour ne pas répéter le thème. On l'utilise aussi pour donner une nouvelle information, lorsque le thème n'est pas encore introduit, par exemple
            \begin{tcolorbox}
                -- 学校の前に、大きい犬がいます。その犬はとてもかわいいです。

                -- がっこう　の　まえ　に、おおきい　いぬ　が　います。その　いぬ　は　とても　かわいい　です。

                -- Devant l'école, il y a un gros chien. C'est un chien très mignon.
            \end{tcolorbox}
            \item \textit{La possession :} On peut exprimer la possession par la structure 
            \begin{center}
                「thème　は　groupe-nominal　が　いる/ある」
            \end{center}
            qui se traduit littéralement par « en ce qui concerne thème, il y a groupe-nominal » et qui signifie « thème a groupe-nominal ».
            \begin{tcolorbox}
                -- 私の部屋には、机とベッドがあります。

                -- わたし　の　へや　には、つくえ　と　ベッド　が　あります。

                -- Dans ma chambre, j'ai un lit et un bureau.
            \end{tcolorbox}
        \end{itemize}
    \end{itemize}

        \subsubsection{から}

    \begin{itemize}
        \item \textbf{Le point de départ :} La particule から indique un point de départ qui peut être temporel ou spacial.
        \item \textbf{Causalité :} Pour exprimer la causalité (« parce que »), on utilise la structure 
        \begin{center}
            「cause　から、conséquence」
        \end{center}
        \begin{tcolorbox}
            -- 私は肉が好きではないから、肉を食べません。

            -- わたし　は　にく　が　すき　ではない　から、にく　を　たべません。

            -- Je ne mange pas de viande car je n'aime pas ça.
        \end{tcolorbox}
    \end{itemize}

        \subsubsection{まで}

    \begin{itemize}
        \item \textbf{Le point d'arrivée :} La particule まで indique un point d'arrivée à la manière de から.
    \end{itemize}

        \subsubsection{ので}

    \begin{itemize}
        \item \textbf{La cause :} La particule ので, à l'instar de la particule から, permet d'exprimer la cause, mais de façon plus douce et plus polie. Elle s'utilise généralement lorsque l'on exprime des faits de façon objective.
    \end{itemize}

        \subsubsection{など}

    \begin{itemize}
        \item \textbf{Liste incomplète :} La particule など, associée à une liste exhaustive de noms liée par や signifie (« etc. »). Associé à un nom seul, elle montre que ce n'est qu'un exemple d'une catégorie plus large.
        
        \begin{tcolorbox}
            -- 雑誌などをよく読みます。パンや卵などを買いました。

            -- ざっしなどを　よく　よみます。ぱんや　たまごやなどをあいました。

            -- Je lis souvent des magazines et des choses comme ça. J'ai acheté du pain, des œufs, et ainsi de suite.
        \end{tcolorbox}
    \end{itemize}

        \subsubsection{くらい}

    \begin{itemize}
        \item \textbf{Approximation :} Lorsque l'on veut exprimer l'idée d'approximation (« environ »), on utilise la particule くらい, qui se lie directement à un nom, verbe ou adjectif.
        
        \begin{tcolorbox}
            -- １０人くらい来ました。３時間くらい勉強しました。

            -- じゅうにんくらい　きました。さんじかんくらい　べんきょうしました。

            -- Environ 10 personnes sont venues. J'ai étudié pendant environ trois heures.
        \end{tcolorbox}
        \item \textbf{Degré ou étendue :} Avec des verbes et des adjectifs, くらい permet d'illustrer à quel point une situation est extrême.
        
        \begin{tcolorbox}
            -- 泣きたいくらい疲れました。涙が出るくらい嬉しいです。

            -- なきたいくらい　つかれました。なみだがでるくらい　うれしい　です。

            -- Je suis tellement fatigué que j'ai envie de pleurer.
        \end{tcolorbox}

        \item \textbf{Strict minimum :} くらい permet d'expliciter ce qui serait le strict minimum, une futilité, ou du moins ce qui pourrait facilement être fait.
            
        \begin{tcolorbox}
            -- 自分の服くらい自分で洗ってください。

            -- じぶんのふくくらい　じぶんで　あらってください。

            -- Lavez au moins vos propres vêtements.
        \end{tcolorbox}
    \end{itemize}

        \subsubsection{より}

    \begin{itemize}
        \item \textbf{Comparaison :} Attaché à un nom, より montre la norme inférieure d'une comparaison (« par rapport à »). À un verbe à la forme du dictionnaire, より établit une comparaison en prenant cette action pour référence.
        
        \begin{tcolorbox}
            -- 今日は昨日より暑いです。食べるより寝るほうが好きです。

            -- きょうは　きのうより　あついです。たべるより　ねる　ほうが　すき　です。

            -- Il fait plus chaud aujourd'hui qu'hier. J'aime plus dormir que manger.
        \end{tcolorbox}
    \end{itemize}

        \subsubsection{のに}

    \begin{itemize}
        \item \textbf{Résultat inattendu :} On utilise のに entre deux propositions lorsque le résultat contredit directement ce à quoi on s'attendrait naturellement d'après le premier fait.
        
        \begin{tcolorbox}
            -- たくさん寝たのに、まだ眠いです

            -- たくさん　ねたのに、まだねむい　です。

            -- Même si j'ai beaucoup dormi, j'ai encore sommeil.
        \end{tcolorbox}

        \item \textbf{Laisser une conclusion non dite :} Dans le langage courant, une phrase peut se terminer complètement par のに pour laisser un sentiment de regret ou une plainte subtile planer dans l'air, laissant l'auditeur deviner la suite.
        
        \begin{tcolorbox}
            -- あんなに頑張ったのに。

            -- あんなに　がんばった　のに。

            -- Même si j'ai fait tellement d'efforts.
        \end{tcolorbox}
    \end{itemize}

    \subsection{Les adjectifs}

    Il existe deux types d'adjectifs en japonais, qui peuvent ou non être conjugués (en temps, en positivité...) à la façon des verbes.

        \subsubsection{Les な-adjectifs}

    Les な-adjectifs sont des adjectifs \textit{invariables}, et c'est le verbe de la phrase qui se conjugue. Les phrases les plus basiques les impliquant impliquent l'attitude d'être, permettant d'user de l'adjectif comme un groupe de mot :
        
    \begin{tcolorbox}
        -- 今日、暇ですか。いいえ、暇じゃないです。

        -- きょう、ひま　です　か。いいえ、ひま　じゃないです。

        -- Es-tu libre aujourd'hui ? Non, je ne le suis pas.
    \end{tcolorbox}

    Leur nom tient du fait que lorsqu'ils précisent un nom, il faut ajouter la particule な entre les deux. 

    \begin{tcolorbox}
        -- きれいな絵でしたか。いいえ、きれいな絵でわありませんでした。

        -- きれい　な　え　でした　か。いいえ、きれい　な　え　でわありませんでした。

        -- Le tableau était-il joli ? Non, ce n'était pas un joli tableau.
    \end{tcolorbox}


        \subsubsection{Les い-adjectifs} 
        
    Les い-adjectifs (dits 形容詞, けいようし) sont des adjectifs \textit{variables}, et on les conjugue d'une façon similaire aux verbes. Pour cela, on repère le radical par l'adjectif sans sa terminaison い, auquel on ajoute des suffixes :
        
    \begin{tblr}[theme = general]{colspec = {Q[c] Q[c] Q[l] Q[l]}}
    & Forme   & Accompli       & Non-accompli \\
    Aff. & Neutre  &\~{}かった       &\~{}い \\
         & Polie   &\~{}かったです    &\~{}いです \\
         & Polie+  &\~{}くありました  &\~{}くあります \\
    Nég. & Neutre  &\~{}くなかった     &\~{}くない \\
         & Polie   &\~{}くなかったです  &\~{}くありません \\
         & Polie+  &\~{}くありませんでした &\~{}くありません \\
    \end{tblr}
        
    \begin{tcolorbox}
        -- その本は面白いですか。いいえ、あまり面白くないです。

        -- その　ほん　は　おもしろい　です　か。いいえ、あまり　おもしろくないです。

        -- Ce livre est-il intéressant ? Non, il n'est pas très intéressant.
    \end{tcolorbox}

    Seule exception notable, l'adjectif いい qui est \textit{irrégulier} dans cette conjugaison : à part sa forme dictionnaire いい, on peut le conjuguer comme on conjuguerait l'adjectif よい (sa forme ancienne). 
        
    Certaines formes verbales (forme en \~{}たい par exemple) s'appliquent aussi aux い-adjectifs, et on peut les conjuguer de la même façon que des verbes.

    \begin{itemize}
        \item \textbf{Les adverbes \~{}く :} Lorsque l'on pose la terminaison\~{}く au い-adjectif, on le transforme en adverbe. Ces adverbes se placent juste avant le verbe pour les caractériser.
        
        \begin{tcolorbox}
            -- 寒くなりました。

            -- さむく　なりました。

            -- Il a commencé à faire froid.
        \end{tcolorbox}

        \item \textbf{Les noms \~{}さ :} Lorsque l'on pose la terminaison\~{}さ au い-adjectif, on le transforme en le nom qui correspond au fait d’être « adjectif ». Par exemple, 速い (はやい) signifie rapide, et 速さ signifie la rapidité.
    \end{itemize}

    \begin{tcolorbox}
        -- 明日は日曜日ですから、私は早く起きたくありません。

        -- あした　は　にちようび　です　から、わたしは　はやく　おきたく　ありません。

        -- Comme demain est dimanche, je ne veux pas me lever tôt.
    \end{tcolorbox}

        \subsubsection{L'utilisation des adjectifs}

    Un adjectif permet souvent à lui seul de donner un sens à la phrase, notamment dans le cas des formulations du type 

    \begin{center}
        「thème　は　sujet　が　adjectif」
    \end{center}

    Des exemples permettent de comprendre la richesse d'une telle formulation :

    \begin{tcolorbox}
        -- 貴方は日本語が上手ですね。いいえ、まだまだです。

        -- あなたは　にほんごが　じょうず　ですね。いいえ、まだまだです。

        -- Tu te débrouilles bien en japonais, n'est-ce pas ? Non, pas encore.
    \end{tcolorbox}

    On utilisera souvent avec cette forme les な-adjectifs 上手 (じょうず) et 下手 (へた), qui signifient respectivement habile, bon et malhabile, mauvais. De la même façon, pour décrire des caractéristiques physiques, nous n'avons pas besoin de verbe, l'adjectif se suffit (mais tout de même conjugué à la forme en て).

    \begin{tcolorbox}
        -- この犬は耳が大きくて、目がかわいいね。

        -- この いぬは みみが おおきくて、めが かわいいね。

        -- Ce chien a de grandes oreilles et de jolis yeux, hein.
    \end{tcolorbox}


    \subsection{Les pronoms}

        \subsubsection{Les préfixes こ\~{} そ\~{} あ\~{} ど\~{}}

    Il s'agit d'un groupe de préfixes qui permettent d'indiquer des choses, lieux et directions en fonction de leur proximité avec le locuteur et l'interlocuteur.

    \begin{tblr}[theme = general]{colspec={Q[l]Q[l]}}
        Préfixe & Signification \\
        こ\~{} & Proche du locuteur \\
        そ\~{} & Proche de l'interlocuteur \\
        あ\~{} & Loin du locuteur et de l'interlocuteur \\
        ど\~{} & Questionne la proximité \\
    \end{tblr}

    On ajoute ensuite un suffixe qui permet d'indiquer différentes choses. 

    \begin{tblr}[theme = general]{colspec={Q[l]Q[l]}}
        Suffixe & Signification \\
       \~{}れ & Indique un objet \\
       \~{}の + nom & Indique un objet identifié \\
       \~{}こ & Indique un lieu \\
       \~{}ちら & Indique une direction ou présente de façon très polie une personne ou un lieu \\
    \end{tblr}

    On pourra ainsi traduire : 
    \begin{itemize}
        \item これ (ceci)、それ (cela)、あれ (cela là-bas)、どれ (lequel)
        \item この (ce + nom-ci)、その (ce nom-là)、あの (ce nom-la là-bas)、どの (quel)
        \item ここ (ici)、そこ (là)、\textcolor{couleurmatiere}{あそこ} (là-bas)、どこ (où)
        \item こちら (vers ici)、そちら (vers là)、あちら (vers là-bas)、どちら (vers où)
    \end{itemize}

    \begin{tcolorbox}
        -- すみません。郵便局は何処ですか。

        -- すみません。ゆうびんきょく　は　どこ　です　か。

        -- Excusez-moi. Où est le bureau de poste ?
    \end{tcolorbox}

    \begin{tcolorbox}
        -- 会社はそちらですか。出口はどちらですか。

        -- こいしゃは　そちら　ですか。でぐちは　どちら　ですか。

        -- Votre entreprise est-elle par là ? Par où est la sortie ?
    \end{tcolorbox}

    
    dare (who), dareno (whose).

        \subsubsection{Le suffixe\~{}んな}

    On peut, avec les préfixes précédemment introduits, aussi utiliser le suffixe\~{}んな, qui aura pour signification principale de décrire (dans la nature ou le type d'un objet) sa proximité avec un objet dont la localisation s'indique par le préfixe.

    Ces démonstratif auront plusieurs fonctions, se rapportant toujours à un objet (caractérisée par sa proximité avec le locuteur ou l'interlocuteur) ou sujet tout juste abordé (par le locuteur ou l'interlocuteur).
    \begin{itemize}
        \item décrire un nom (dans son type ou sa nature) \textit{semblable à quelque chose  physiquement caractérisé par sa proximité}.
        \begin{tcolorbox}
            -- こんなカメラが欲しいです。そんな鞄はどこで買いましたか。あんな山に登りたいです。

            -- こんな　カメラが　ほしい　です。そんな　かばんは　どこで　あいましたか。あんな　やまに　のぼりたいです。

            -- Je veux ce genre d'appareil photo. Où avez-vous acheté ce genre de sac ? Je veux grimper sur ce genre de montagne.
        \end{tcolorbox}
        \item exprimer \textit{la surprise ou l'émotion}. Avec そんな, par rapport à un sujet tout juste introduit par l'interlocuteur. Avec あんな, par rapport à des souvenirs communs (pas nécessairement avec émotion).
        \begin{tcolorbox}
            -- こんな所にレストランがあります。そんなことを言わないでください。あんな話は嘘です。

            -- こんな　ところに　レストランが　あります。そんな　ことを　いわないでください。あんな　はなしは　うそ　です。

            -- Il y a un restaurant dans un endroit comme ça !? Ne dites pas une telle chose s'il vous plaît. Ce genre d'histoires sont des mensonges.
        \end{tcolorbox}
        \item montrer \textit{l'aversion ou le rejet} avec こんな.
        \begin{tcolorbox}
            -- んな雑誌は読みません。

            -- こんな　ざっしは　よみません。

            -- Je ne lis pas ce genre de magazine.
        \end{tcolorbox}
    \end{itemize}
    Le mot interrogatif どんな permet principalement de demander une description sur le nom qu'il précède ou le caractère (personnalité, apparence) d'une personne.
    \begin{tcolorbox}
        -- どんな音楽が好きですか。新しい先生はどんな人ですか。

        -- どんな おんがくが　すき　ですか。あたらしい　せんせいは　どんな　ひと　ですか。

        -- Quel genre de musique aimez-vous ? Quel genre de personne est le nouveau professeur ?
    \end{tcolorbox}

        \subsubsection{Mots interrogatifs}

    \begin{itemize}
        \item 誰 (だれ) est un pronom interrogatif signifiant « qui ». On peut l'utiliser pour demander l'identité d'une personne ou l'appartenance d'un objet.
        
        \begin{tcolorbox}
        -- これは誰の鞄ですか。週末、誰と遊びに行きますか。

        -- これ　は　だれ　の　かばん　です　か。しゅうまつ、だれと　あそびにいきますか。

        -- À qui est ce sac-ci ? Avec qui allez-vous sortir ce week-end ?
        \end{tcolorbox}

        Dans une situation formelle, on utilisera plutôt どなた.

        \item Pour demander \textit{une impression ou une opinion}, \textit{un état ou une condition}, \textit{des méthodes ou des choix}, on utilise le mot interrogatif どう, qui peut souvent se traduire par « comment » ou « que ». Il se place, à la façon d'un adverbe, avant le verbe.
        
        \begin{tcolorbox}
            -- 日本の生活はどうですか。今日の天気はどうですか。此れからどうしますか。

            -- にほんの　せいかつは　どう　ですか。きょうの　てんきは　どう　ですか。これから　どう　しますか。

            -- Comment est la vie au Japon ? Quel temps fait-il aujourd'hui ? Qu'allez-vous faire désormais ?
        \end{tcolorbox}

        Son équivalent poli est le mot interrogatif いかが, qui sert à interroger de façon respecteuse sur un désir, une opinion ou un état.

        \begin{tcolorbox}
            -- お茶はいかがですか。

            -- おちゃは　いかが　ですか。

            -- Souhaitez-vous du thé ?
        \end{tcolorbox}

        \item 何 (なん, なに) est un mot interrogatif fondamental, qui permet de demander l'identité d'un objet, d'une action ou d'une situation, souvent traduit par « que » ou « quoi ». Son utilisation est très large, et dépend majoritairement du bon emploi de la particule qui le suis. Il faudra donc connaître le sens de ses liens aux particules.
        
        \begin{tcolorbox}
            -- なにから始めますか。何もありません。

            -- なにから はじめますか。なにもありません。

            -- Par quoi allez-vous commencer ? Il n'y a rien.
        \end{tcolorbox}

        \item どうして est un mot interrogatif courant qui permet de demander des raisons (« pourquoi »).
        
        \begin{tcolorbox}
            -- どうして日本語を勉強しますか。えっ、どうして知っているんですか。

            -- どうして　にほんごを　べんきょうしますか。えっ、どうして　しっているんてすか。

            -- Pourquoi étudiez-vous le japonais? Attendez, pourquoi savez-vous cela?
        \end{tcolorbox}

        De façon beaucoup plus familière, on utilisera le mot interrogatif なんで, pour demander des raisons, souvent avec un léger ton de surprise ou de curiosité.

        \begin{tcolorbox}
            -- なんで行かないの？

            -- なんで　いかないの?

            -- Pourquoi tu n'y vas pas?
        \end{tcolorbox}

        \item 幾ら (いくら) est le mot interrogatif pour demander un prix.
    \end{itemize}


    \subsection{Les adverbes}

    de fréquence, à utiliser avec un verbe à l'affirmatif.
    \begin{itemize}
        \item 毎日 (まいにち, tous les jours) 
        \item よく (souvent)
        \item ときどき (parfois)
    \end{itemize}
    de fréquence, à utiliser avec un verbe au négatif.
    \begin{itemize}
        \item ぜんぜん (jamais, pas du tout)
        \item あまり (pas souvent, pas beaucoup)
        
        たけしさんはあまり勉強しません。

        たけしさん　は　あまり　べんきょうしません。

        Takeshi n'étudie pas beaucoup.

    \end{itemize}

    \subsection{Les compteurs}

    Les classificateurs numéraux permettent d'exprimer le décompte d'un certain type d'objet ou être vivant. Chaque classificateur (aussi appelé compteur) est associé à une catégorie de noms bien définie : 
    
    \begin{tblr}[theme = general]{colspec={Q[l,1]Q[l,2]Q[l,6]}}
        Compteur & 仮名 & Comptés \\
        人 & ニン & êtres humains \\
        匹 & ヒキ & petits animaux \\
        頭 & トウ & grands animaux \\
        羽 & ワ & animaux ailés (et lapins) \\
        枚 & マイ & objets plats et fins \\
        本 & ホン & objets longs et cylindriques \\
        冊 & サツ & objets reliés \\
        台 & ダイ & machines, produits industriels, véhicules \\
        杯 & ハイ & contenu d'un verre, d'un bol, d'une cuillère \\
        個 & コ & choses de petite et moyenne tailles, de formes variées \\
        つ & つ & choses concrètes ou abstraites, difficiles à classer
    \end{tblr}

    Certaines prononciations varient avec la combinaison des nombres et des classeurs, non détaillés ici. Retenons simplement que les classificateurs notés en 片仮名 s'utilisent avec les prononciations courantes des nombres (sauf 一人 [ひとり]  et 二人 [ふたり]), et les classificateurs notés en 平仮名 avec leur prononciation purement japonaise (つ).

    On doit systématiquement préciser la catégorie d'objets que l'on compte, et on place cet indicateur juste avant le verbe ou à l'aide de la particule の.

    \begin{tcolorbox}
        -- 猫が\textcolor{couleurmatiere}{二匹}いる。/　\textcolor{couleurmatiere}{二匹}\textcolor{gray}{の}猫がいる。

        -- ねき　が　にひき　いる。/　にひき　の　ねこ　が　いる。

        -- Il y a deux chats. / Il y a les deux chats.
    \end{tcolorbox}

    Pour poser des questions de quantité, on utilise 何 suivi du classificateur, sauf pour つ qui se demandera par いくつ.

    \begin{tcolorbox}
        -- 空には鳥が\textcolor{couleurmatiere}{何羽}いるか。/　空には\textcolor{couleurmatiere}{何羽}\textcolor{gray}{の}鳥がいるか。

        -- そら　には　とり　が　なんわ　いる　か。/　そら　には　なんわ　の　とり　が　いる　か。

        -- Combien y'a-t-il d'oiseaux dans le ciel ? / Combien d'oiseaux y'a-t-il dans le ciel ?
    \end{tcolorbox}

        \subsubsection{Les nombres ordinaux}

    Il existe deux façons de créer des nombres ordinaux :
    \begin{itemize}
        \item \textbf{Le préfixe 第\~ :} En ajoutant le préfixe 第\~{} devant un nombre cardinal, on exprime sa version ordinale qui peut être utilisée uniquement dans un contexte général prédéfini. Par exemple, dans une course de chevaux où ceux-ci sont numérotés,
        \begin{tcolorbox}
            -- \textcolor{couleurmatiere}{第三}\textcolor{gray}{の}馬に賭けました。

            -- だいさん　の　うま　に　かけました。

            -- J'ai parié sur le troisième cheval.
        \end{tcolorbox}
        \item \textbf{Le suffixe\~{}目 :} En ajoutant le suffixe\~{}目 derrière un nombre associé à un compteur, on exprime sa version ordinale qui peut être utilisée dans tous les cas. 
        \begin{tcolorbox}
            -- 右から\textcolor{couleurmatiere}{三頭目}\textcolor{gray}{の}馬が一番きれいです。

            -- みぎ　から　さんとうめ　の　うま　が　いちばん　きれい　です。

            -- Le troisième cheval en partant de la droite est le plus beau.
        \end{tcolorbox}
    \end{itemize}
    
